% Пример на русском языке: короткая заметка, в которой порог поставлен вручную.
% Сборка: lualatex report-ru.tex, два раза. Шрифты лежат в ../doc/fonts.
\documentclass[11pt,a4paper]{article}
\usepackage[a4paper,left=58mm,right=17mm,top=30mm,bottom=29mm,headsep=8mm,footskip=11mm]{geometry}
\usepackage{fontspec}
\setmainfont{Lora}[Path=../doc/fonts/,Extension=.ttf,UprightFont=*-Regular,BoldFont=*-Bold,
  ItalicFont=*-Italic,BoldItalicFont=*-BoldItalic]
\setsansfont{Inter}[Path=../doc/fonts/,Extension=.otf,UprightFont=*-Regular,BoldFont=*-SemiBold,
  ItalicFont=*-Italic,BoldItalicFont=*-Bold,Scale=0.9]
\linespread{1.1}
% Если в установке TeX нет русских переносов, берём их из ../doc (hyph-utf8, LPPL)
\makeatletter
\ifdefined\l@russian\else
  \newlanguage\l@russian
  \directlua{
    local function slurp(n) local f = assert(io.open(n, "r")) local s = f:read("*a") f:close() return s end
    local l = lang.new(\the\l@russian)
    lang.patterns(l, slurp("../doc/hyph-ru.pat.txt"))
    lang.hyphenation(l, slurp("../doc/hyph-ru.hyp.txt"))}
\fi
\makeatother
\usepackage[provide=*]{babel}
\babelprovide[import,main,hyphenrules=russian]{russian}
\lefthyphenmin=2 \righthyphenmin=2
\usepackage[palette=ultramarine]{limen}
\usepackage{fancyhdr}
\pagestyle{fancy}\fancyhf{}\renewcommand\headrulewidth{0pt}
\fancyhead[R]{\footnotesize\sffamily\color{gray}Как лазер переходит порог}
\fancyfoot[R]{\footnotesize\sffamily\color{gray}\thepage}
\makeatletter
\renewcommand\section{\@startsection{section}{1}{\z@}{-2.6ex plus -.6ex minus -.2ex}{1.2ex plus .2ex}%
  {\normalfont\large\bfseries\color{limenink}}}
\makeatother
\begin{document}
\noindent{\LARGE\bfseries\color{limenink}Как лазер переходит порог\par}
\medskip
\noindent{\itshape\color{gray}Заметка-пример к пакету limen\par}
\bigskip

\limeninitial{Л}азер устроен просто: это усилитель света, который замкнули сам на себя с помощью двух зеркал. Кристалл
или газ внутри усиливает свет, зеркала возвращают его обратно, и каждый новый проход начинается с того, чем закончился
предыдущий. Самое интересное в таком устройстве происходит не плавно, а в одной точке. Пока накачка слаба, лазер почти
не светит. Стоит ей перейти определённое значение, и он начинает работать совсем иначе. Эту точку называют порогом.

На полях этой заметки происходит то же самое. Пока идёт рассказ о том, что бывает ниже порога, луч остаётся тонким
и почти белым, а поле тёмным. Там, где в тексте появляется порог, начинается рост.

\section{Ниже порога}

Накачка переводит атомы рабочего вещества на верхний уровень. В лазере на алюмоиттриевом гранате с неодимом такой
уровень живёт долго, около 230~мкс, поэтому атомы успевают накопиться. Их избыток над нижним уровнем называют инверсной
населённостью. Именно он даёт усиление: фотон подходящей частоты, пролетая мимо возбуждённого атома, вынуждает его
отдать энергию и порождает второй такой же фотон.

Но зеркала не идеальны, и часть света теряется на каждом обходе. Часть теряется намеренно, ведь одно из зеркал должно
выпускать луч наружу. Пока усиление за обход меньше потерь, любой фотон в резонаторе быстро исчезает. Кристалл при этом
светится, но это обычное спонтанное излучение: слабое, разбегающееся во все стороны и широкое по спектру. Если увеличить
накачку вдвое, свечения станет примерно вдвое больше, и ничего принципиально нового не произойдёт.

Ниже порога лазер ведёт себя как лампа, причём довольно плохая. Почти вся энергия накачки уходит в тепло и в свечение,
которое никто не собирает. Судить по этому режиму о том, на что способно устройство, было бы ошибкой.

\section{Порог}
\limenthreshold

Порог наступает в тот момент, когда усиление за полный обход резонатора сравнивается с потерями за обход. С этого момента
фотон, идущий вдоль оси, в среднем не исчезает: на его место приходит новый. Число фотонов начинает расти лавинообразно,
и рост останавливается только тогда, когда сами фотоны заметно снимают инверсию.

Здесь и появляется главная особенность лазера. Усиление не может подняться выше потерь надолго, иначе число фотонов росло
бы без конца. Поэтому выше порога инверсия почти перестаёт меняться. Её значение как будто зажато на пороговом уровне,
сколько бы энергии ни подавали. Вся добавка накачки уходит уже не в запас атомов, а в свет.

\limenrule

\section{Выше порога}

Теперь выходная мощность растёт почти линейно с накачкой: $P \approx \eta\,(P_\mathrm{p} - P_\mathrm{th})$. Наклон
этой прямой называют дифференциальной эффективностью. Он не может быть больше отношения длин волн накачки и генерации,
потому что один фотон накачки даёт не больше одного фотона лазера. Для накачки на 808~нм и излучения на 1064~нм этот
предел составляет около 76~\%.

Меняется и сам свет. Он идёт узким пучком вдоль оси резонатора, потому что всё остальное гибнет на зеркалах. Линия
в спектре резко сужается, потому что лавину выигрывают те частоты, для которых усиление чуть выше. Волны складываются
в фазе, и свет становится когерентным. Всё это одна и та же лавина, увиденная с разных сторон.

На полях веретёна становятся шире, а поле набирает цвет. Это уже не физика, а образ: настоящий луч неодимового лазера
одноцветный и невидимый глазу, его длина волны 1064~нм лежит в инфракрасной области. Радугу \limenspectrum[2.5cm] он не
образует, а звёздочка \limenspark{} на его пути ничего не значит.

\section{Откуда всё-таки может взяться радуга}

Сделать из лазерного луча белый свет можно, и это обычная лабораторная техника. Короткие мощные импульсы направляют
в тонкое нелинейное волокно. Там спектр расширяется настолько, что перекрывает всю видимую область. Такое излучение
называют суперконтинуумом. Если пропустить его через призму, получится именно та картина, что на последней странице.

Так что на полях нарисован не неодимовый лазер, а история о нём. Сначала почти ничего, потом порог, потом рост и в конце
свет, который уходит за край листа.
\end{document}
